\documentclass{ximera}

\input{../../../preamble.tex}


\definecolor{fvdnavy}{HTML}{071D43}
\definecolor{fvdcyan}{HTML}{20BFC6}
\definecolor{fvdcyanlight}{HTML}{EFFCFB}
\definecolor{fvdmagenta}{HTML}{F10D69}
\definecolor{fvdmagentalight}{HTML}{FFF1F7}
\definecolor{fvdtext}{HTML}{163044}





%=================================================
% Pedagogische omgevingen
% PDF: tcolorbox
% HTML: learning-box
%=================================================

\newenvironment{voorkennis}
{%
    \ifdefined\HCode
        \HCode{%
<section class="learning-box learning-box--prior">
<div class="learning-box__header">
<span class="learning-box__icon" aria-hidden="true"></span>
<span class="learning-box__title">Voorkennis</span>
</div>
<div class="learning-box__content">
        }%
        \begin{itemize}
    \else
       \begin{tcolorbox}[
    enhanced,
    breakable,
    colback=fvdcyanlight,
    colframe=fvdcyan,
    coltext=fvdtext,
    boxrule=0.5pt,
    leftrule=4pt,
    arc=4mm,
    outer arc=4mm,
    left=7mm,
    right=6mm,
    top=5mm,
    bottom=5mm,
    before skip=6mm,
    after skip=6mm,
    title=Voorkennis,
    coltitle=fvdcyan,
    fonttitle=\bfseries\Large,
    attach title to upper={\par\smallskip},
    boxed title style={
        empty,
        left=0mm,
        right=0mm,
        top=0mm,
        bottom=0mm
    },
    overlay={
        \draw[
            fvdcyan,
            line width=0.5pt
        ]
        ([xshift=42mm,yshift=-13mm]frame.north west)
        --
        ([xshift=-7mm,yshift=-13mm]frame.north east);
    }
]
        \begin{itemize}
    \fi
}
{%
    \end{itemize}
    \ifdefined\HCode
        \HCode{%
</div>
</section>
        }%
    \else
        \end{tcolorbox}
    \fi
}


\newenvironment{leerdoelen}
{%
    \ifdefined\HCode
        \HCode{%
<section class="learning-box learning-box--goals">
<div class="learning-box__header">
<span class="learning-box__icon" aria-hidden="true"></span>
<span class="learning-box__title">Leerdoelen</span>
</div>
<div class="learning-box__content">
        }%
        \begin{itemize}
    \else
\begin{tcolorbox}[
    enhanced,
    breakable,
    colback=fvdmagentalight,
    colframe=fvdmagenta,
    coltext=fvdtext,
    boxrule=0.5pt,
    leftrule=4pt,
    arc=4mm,
    outer arc=4mm,
    left=7mm,
    right=6mm,
    top=5mm,
    bottom=5mm,
    before skip=6mm,
    after skip=6mm,
    title=Leerdoelen,
    coltitle=fvdmagenta,
    fonttitle=\bfseries\Large,
    attach title to upper={\par\smallskip},
    boxed title style={
        empty,
        left=0mm,
        right=0mm,
        top=0mm,
        bottom=0mm
    },
    overlay={
        \draw[
            fvdmagenta,
            line width=0.5pt
        ]
        ([xshift=42mm,yshift=-13mm]frame.north west)
        --
        ([xshift=-7mm,yshift=-13mm]frame.north east);
    }
]
        \begin{itemize}
    \fi
}
{%
    \end{itemize}
    \ifdefined\HCode
        \HCode{%
</div>
</section>
        }%
    \else
        \end{tcolorbox}
    \fi
}


\addPrintStyle{../../..}

\title{Samenhang is nog geen oorzaak}
\author{Ingmar Herreman}

\begin{abstract}
Een puntenwolk kan een duidelijk patroon vertonen. De
correlatiecoëfficiënt kan dicht bij $1$ of $-1$ liggen en een
regressiemodel kan de gegevens uitstekend beschrijven.

Maar daarmee is nog niet aangetoond dat de ene grootheid de andere
\textbf{veroorzaakt}.

In dit hoofdstuk onderzoeken we het verschil tussen
\textbf{samenhang} en \textbf{causaliteit}. We bekijken hoe een
waargenomen samenhang kan ontstaan door een gemeenschappelijke derde
variabele, door omgekeerde beïnvloeding, door toeval of door de manier
waarop gegevens verzameld werden.

We leren daarom niet alleen vragen
\emph{``Is er een verband?''}, maar ook
\emph{``Welke conclusie laten deze gegevens werkelijk toe?''}

Zo maken we de volledige beweging van werkelijkheid naar data en
wiskundig model --- en opnieuw terug naar de werkelijkheid.
\end{abstract}

\begin{document}

% FVD_CHAPTER_DATA_START
% Automatisch gegenereerd uit index.tex.
% Niet handmatig aanpassen.
\fvdchapterdata{3}{Verbanden onderzoeken met data}{10}
% FVD_CHAPTER_DATA_END



\maketitle

\begin{voorkennis}
  \item Je kan een spreidingsdiagram analyseren.
  \item Je kan de correlatiecoëfficiënt $r$ interpreteren.
  \item Je weet dat $r$ de richting en sterkte van een lineair verband beschrijft.
  \item Je kan een lineair regressiemodel opstellen en interpreteren.
  \item Je kan interpolatie en extrapolatie onderscheiden.
  \item Je kan een wiskundig model terugkoppelen naar de context.
\end{voorkennis}

\begin{leerdoelen}
  \item Je kan in concrete situaties het verschil uitleggen tussen samenhang en causaliteit.
  \item Je kan aangeven welke conclusies een correlatie of regressiemodel wel en niet ondersteunt.
  \item Je kan herkennen dat een derde variabele een waargenomen samenhang kan helpen verklaren.
  \item Je kan omgekeerde causaliteit als mogelijke verklaring herkennen.
  \item Je kan uitleggen waarom observationele gegevens meestal onvoldoende zijn om causaliteit vast te stellen.
  \item Je kan uitleggen waarom een gecontroleerd experiment sterkere causale uitspraken kan ondersteunen.
  \item Je kan een causale claim kritisch beoordelen op basis van gegevens en onderzoeksopzet.
  \item Je kan een voorzichtige, wetenschappelijk verantwoorde conclusie formuleren.
\end{leerdoelen}


% ============================================================
\FVDsection{Een sterk verband --- wat weten we nu?}
% ============================================================

Stel dat we twee numerieke grootheden $x$ en $y$ onderzoeken.

We maken een spreidingsdiagram en vinden een bijna rechte puntenwolk.
ICT geeft bovendien

\[
r=0{,}96.
\]

Een lineaire regressie levert

\[
\hat y=2{,}4x+7{,}1.
\]

We mogen dan besluiten dat er in deze dataset een
\textbf{sterke positieve lineaire samenhang} bestaat.

Maar mogen we ook besluiten:

\begin{center}
\textbf{``Een stijging van $x$ veroorzaakt een stijging van $y$''?}
\end{center}

Nee. Daarvoor hebben we meer informatie nodig.

\begin{opmerking}
Correlatie en regressie beschrijven een patroon in gegevens.

Ze beantwoorden vragen zoals:

\begin{itemize}
  \item veranderen twee grootheden samen?
  \item in welke richting?
  \item hoe sterk is hun lineaire samenhang?
  \item welk lineair model beschrijft het patroon?
\end{itemize}

Ze vertellen op zichzelf niet \textbf{waarom} dat patroon bestaat.
\end{opmerking}

Daarmee komen we bij een fundamenteel onderscheid.


% ============================================================
\FVDsection{Samenhang en causaliteit}
% ============================================================

\begin{definitie}
We spreken van \textbf{samenhang} wanneer waarden van twee grootheden
op een systematische manier samen voorkomen of samen veranderen.
\end{definitie}

Een correlatiecoëfficiënt is een maat voor een bepaalde vorm van
samenhang, namelijk \textbf{lineaire samenhang}.

\begin{definitie}
We spreken van \textbf{causaliteit} wanneer een verandering in een
grootheid een verandering in een andere grootheid
\textbf{veroorzaakt}.
\end{definitie}

Dat zijn verschillende uitspraken.

\[
\boxed{
\text{samenhang}
\neq
\text{causaliteit}
}
\]

\begin{voorbeeld}[Een wetenschappelijke formulering]
Stel dat leerlingen die meer uren slapen gemiddeld betere scores halen
op een concentratietest.

Uit de gegevens kunnen we bijvoorbeeld besluiten:

\begin{quote}
In de onderzochte groep bestaat een positieve samenhang tussen
slaapduur en de score op de concentratietest.
\end{quote}

Maar uit die samenhang alleen volgt nog niet:

\begin{quote}
Meer slapen veroorzaakt noodzakelijk een hogere score.
\end{quote}

Misschien verschillen de groepen ook in stress, gezondheid,
schermgebruik, leeftijd of andere factoren.
\end{voorbeeld}


% ============================================================
\FVDsection{Waarom kan samenhang ontstaan zonder directe oorzaak?}
% ============================================================

Wanneer twee grootheden samenhangen, zijn meerdere verklaringen
mogelijk.

Een goede onderzoeker probeert die mogelijkheden van elkaar te
onderscheiden.


% ------------------------------------------------------------
\FVDsubsection{Een derde variabele}
% ------------------------------------------------------------

Soms worden twee grootheden allebei beïnvloed door een derde grootheid.

\begin{voorbeeld}[IJsverkoop en verdrinkingsincidenten]
Stel dat in een dataset op dagen met een hogere ijsverkoop ook meer
verdrinkingsincidenten voorkomen.

We zouden een positieve correlatie kunnen vinden tussen

\[
\text{ijsverkoop}
\quad\text{en}\quad
\text{aantal verdrinkingsincidenten}.
\]

Daaruit volgt uiteraard niet dat ijs eten verdrinkingen veroorzaakt.

Een mogelijke derde variabele is de \textbf{temperatuur}:

\[
\text{hogere temperatuur}
\longrightarrow
\begin{cases}
\text{meer ijsverkoop},\\
\text{meer mensen die gaan zwemmen}.
\end{cases}
\]

De temperatuur kan dus bijdragen aan de samenhang tussen beide
gemeten grootheden.
\end{voorbeeld}

\begin{definitie}
Een \textbf{derde variabele} is een bijkomende variabele die de
waargenomen samenhang tussen twee onderzochte grootheden mee kan
verklaren.
\end{definitie}

\begin{opmerking}
Het noemen van een mogelijke derde variabele bewijst nog niet dat die
variabele werkelijk de volledige verklaring is.

Het toont wel aan waarom de oorspronkelijke correlatie op zichzelf
geen bewijs voor een directe causale relatie vormt.
\end{opmerking}


% ------------------------------------------------------------
\FVDsubsection{Omgekeerde causaliteit}
% ------------------------------------------------------------

Zelfs als twee grootheden causaal met elkaar verbonden zijn, is de
richting niet altijd meteen duidelijk.

\begin{voorbeeld}[Beweging en lichamelijke conditie]
Stel dat mensen die meer bewegen gemiddeld een betere lichamelijke
conditie hebben.

Een mogelijke verklaring is

\[
\text{meer bewegen}
\longrightarrow
\text{betere conditie}.
\]

Maar ook de omgekeerde richting kan een rol spelen:

\[
\text{betere conditie}
\longrightarrow
\text{meer kunnen of willen bewegen}.
\]

In werkelijkheid kunnen beide effecten tegelijk aanwezig zijn.
\end{voorbeeld}

\begin{definitie}
Bij \textbf{omgekeerde causaliteit} is het mogelijk dat de richting
van oorzaak en gevolg anders is dan eerst wordt verondersteld.
\end{definitie}


% ------------------------------------------------------------
\FVDsubsection{Toeval}
% ------------------------------------------------------------

Wanneer we veel variabelen onderzoeken, kunnen sommige toevallig
samenhangen.

Een correlatie in één dataset is dus niet automatisch een universeel
patroon.

Bij kleine datasets is voorzichtigheid extra belangrijk: enkele
waarnemingen kunnen een groot effect hebben op de gevonden samenhang.

\begin{opmerking}
Een opvallende correlatie is een reden om verder te onderzoeken,
niet automatisch een bewijs van een mechanisme.
\end{opmerking}


% ------------------------------------------------------------
\FVDsubsection{De manier waarop gegevens verzameld worden}
% ------------------------------------------------------------

Ook de selectie van waarnemingen kan het beeld beïnvloeden.

Stel dat we het verband tussen beweging en gezondheid willen
onderzoeken, maar alleen gegevens verzamelen bij leden van een
sportclub.

Dan onderzoeken we geen willekeurige doorsnede van de volledige
bevolking.

De manier waarop een steekproef tot stand komt, kan dus mee bepalen
welke patronen we zien.

\begin{opmerking}
Een statistisch patroon moet altijd gelezen worden samen met de vraag:

\[
\important{
\text{Waar komen deze gegevens vandaan?}
}
\]
\end{opmerking}


% ============================================================
\FVDsection{Observationeel onderzoek}
% ============================================================

Veel datasets ontstaan doordat onderzoekers de werkelijkheid
\textbf{observeren} zonder zelf een grootheid toe te wijzen of te
veranderen.

\begin{definitie}
Bij een \textbf{observationeel onderzoek} observeert en meet de
onderzoeker bestaande verschillen zonder de onderzochte situatie
doelbewust toe te wijzen.
\end{definitie}

Voorbeelden zijn:

\begin{itemize}
  \item slaapduur en schoolresultaten bij leerlingen meten;
  \item luchtvervuiling en gezondheidsgegevens in verschillende steden vergelijken;
  \item temperatuur en energieverbruik van een gebouw registreren;
  \item beweging en conditie bij een groep personen meten.
\end{itemize}

Observationeel onderzoek kan zeer waardevol zijn. Het kan patronen
zichtbaar maken, hypotheses opleveren en voorspellingen ondersteunen.

Maar mogelijke derde variabelen zijn vaak moeilijk volledig uit te
sluiten.

Daarom zijn causale uitspraken op basis van uitsluitend observationele
samenhang meestal niet gerechtvaardigd.


% ============================================================
\FVDsection{Experimenten en causaliteit}
% ============================================================

Wanneer we een mogelijke oorzaak werkelijk willen onderzoeken, proberen
we de onderzoeksopzet zo te organiseren dat alternatieve verklaringen
zoveel mogelijk worden beperkt.

\begin{definitie}
Bij een \textbf{experiment} grijpt de onderzoeker doelbewust in en
onderzoekt hij wat er gebeurt wanneer een bepaalde factor wordt
veranderd, terwijl andere relevante omstandigheden zoveel mogelijk
worden gecontroleerd.
\end{definitie}

\begin{voorbeeld}[Meststof en plantengroei]
We willen onderzoeken of een bepaalde hoeveelheid meststof de groei van
planten beïnvloedt.

Alleen planten observeren die toevallig verschillende hoeveelheden
meststof krijgen, laat veel alternatieve verklaringen toe.

Een sterker ontwerp is bijvoorbeeld:

\begin{itemize}
  \item vergelijkbare planten gebruiken;
  \item planten willekeurig over behandelingsgroepen verdelen;
  \item de hoeveelheid meststof doelbewust variëren;
  \item licht, water, temperatuur en andere omstandigheden zo gelijk mogelijk houden;
  \item daarna de groei vergelijken.
\end{itemize}

Zo wordt het aannemelijker dat een systematisch verschil in groei
samenhangt met de onderzochte behandeling.
\end{voorbeeld}


% ------------------------------------------------------------
\FVDsubsection{Waarom willekeurig verdelen?}
% ------------------------------------------------------------

Wanneer proefeenheden willekeurig over groepen worden verdeeld,
proberen we te vermijden dat één groep vooraf systematisch anders is
dan een andere.

\begin{opmerking}
Randomisatie maakt een experiment niet automatisch perfect.

Ze helpt wel om alternatieve verklaringen te beperken en versterkt
daardoor de mogelijkheid om een causaal effect te onderzoeken.
\end{opmerking}


% ------------------------------------------------------------
\FVDsubsection{Waarom controleren?}
% ------------------------------------------------------------

Als verschillende factoren tegelijk veranderen, weten we niet welke
factor verantwoordelijk is voor een waargenomen verschil.

Daarom proberen we in een gecontroleerd experiment relevante
omstandigheden gelijk te houden.

\[
\important{
\text{één factor onderzoeken}
\quad\Rightarrow\quad
\text{andere relevante factoren zo goed mogelijk controleren}
}
\]


% ============================================================
\FVDsection{Wat mag je uit een onderzoek besluiten?}
% ============================================================

We kunnen nu verschillende soorten uitspraken onderscheiden.

\begin{center}
\begin{tabular}{p{0.30\textwidth}|p{0.60\textwidth}}
\textbf{Informatie} & \textbf{Wat ondersteunt ze?}\\ \hline
Spreidingsdiagram &
Een visuele beschrijving van samenhang in de waargenomen gegevens.\\
Correlatiecoëfficiënt &
Richting en sterkte van lineaire samenhang.\\
Regressiemodel &
Een wiskundige beschrijving en mogelijke voorspellingen binnen een geschikt toepassingsgebied.\\
Observationeel onderzoek &
Samenhang vaststellen; causaliteit blijft doorgaans onzeker.\\
Goed gecontroleerd experiment &
Kan, afhankelijk van de onderzoeksopzet, sterkere causale conclusies ondersteunen.
\end{tabular}
\end{center}

\begin{opmerking}
Het woord \textbf{veroorzaakt} vraagt dus veel meer onderbouwing dan
woorden zoals

\begin{itemize}
  \item hangt samen met;
  \item gaat samen met;
  \item is geassocieerd met;
  \item wordt door het model voorspeld.
\end{itemize}
\end{opmerking}


% ============================================================
\FVDsection{Van krantenkop naar wetenschappelijke uitspraak}
% ============================================================

In nieuwsberichten en sociale media wordt een statistische samenhang
soms vertaald naar een causale titel.

Bijvoorbeeld:

\begin{quote}
``Leerlingen die ontbijten halen hogere cijfers: ontbijt maakt je
slimmer.''
\end{quote}

Stel dat het onderliggende onderzoek alleen heeft vastgesteld dat
leerlingen die regelmatig ontbijten gemiddeld hogere cijfers behalen.

Dan ondersteunt het onderzoek voorlopig slechts een uitspraak over
\textbf{samenhang}.

Andere factoren kunnen een rol spelen, bijvoorbeeld slaapritme,
thuissituatie, gezondheid of algemene leefgewoonten.

Een wetenschappelijk voorzichtiger formulering is:

\begin{quote}
In de onderzochte groep hangt regelmatig ontbijten samen met hogere
schoolresultaten.
\end{quote}

\begin{opmerking}
Kritisch omgaan met gegevens betekent niet dat we elke mogelijke
causale verklaring afwijzen.

Het betekent dat we het onderscheid bewaken tussen

\[
\important{
\text{wat de gegevens suggereren}
\quad\text{en}\quad
\text{wat het onderzoek werkelijk aantoont}.
}
\]
\end{opmerking}


% ============================================================
\FVDsection{Een vaste strategie voor causale claims}
% ============================================================

Wanneer iemand uit gegevens een oorzaak-gevolgrelatie afleidt, stel je
achtereenvolgens deze vragen:

\begin{enumerate}
  \item \textbf{Wat is precies waargenomen?}\\
        Welke twee grootheden vertonen samenhang?

  \item \textbf{Hoe sterk en van welke vorm is de samenhang?}\\
        Wat tonen puntenwolk, correlatie en eventueel regressie?

  \item \textbf{Hoe zijn de gegevens verzameld?}\\
        Observationeel of experimenteel?

  \item \textbf{Kan een derde variabele een rol spelen?}\\
        Welke alternatieve verklaringen zijn mogelijk?

  \item \textbf{Kan de richting omgekeerd zijn?}\\
        Is het duidelijk welke variabele oorzaak en welke gevolg is?

  \item \textbf{Zijn relevante factoren gecontroleerd?}\\
        Hoe sterk is de onderzoeksopzet?

  \item \textbf{Welke conclusie is werkelijk verantwoord?}\\
        Spreek over samenhang als causaliteit niet voldoende is aangetoond.
\end{enumerate}

Deze strategie is breder bruikbaar dan statistiek. Ze is een manier
van wetenschappelijk redeneren.


% ============================================================
\FVDsection{De volledige leerlijn van dit thema}
% ============================================================

We begonnen met ruwe gegevens.

Daaruit ontstond een puntenwolk:

\[
(x_i,y_i)
\longrightarrow
\text{spreidingsdiagram}.
\]

Vervolgens beschreven we het patroon:

\[
\text{richting}
+
\text{vorm}
+
\text{sterkte}
+
\text{bijzonderheden}.
\]

Voor ongeveer lineaire verbanden gebruikten we de
correlatiecoëfficiënt:

\[
-1\leq r\leq1.
\]

Daarna bouwden we een lineair model:

\[
\hat y=ax+b.
\]

Met dat model konden we voorspellen en het verschil tussen
interpolatie en extrapolatie onderzoeken.

Maar pas daarna komt de wetenschappelijk beslissende stap:

\[
\important{
\text{Wat betekent dit resultaat in de werkelijkheid?}
}
\]

De volledige cyclus wordt dus

\[
\boxed{
\begin{array}{c}
\text{werkelijkheid}\\
\downarrow\\
\text{gegevens}\\
\downarrow\\
\text{puntenwolk}\\
\downarrow\\
\text{correlatie}\\
\downarrow\\
\text{regressiemodel}\\
\downarrow\\
\text{voorspelling en conclusie}\\
\downarrow\\
\text{kritische terugkeer naar de werkelijkheid}
\end{array}
}
\]

Dat laatste is essentieel.

Wiskunde helpt ons patronen zichtbaar te maken en modellen te bouwen,
maar de betekenis van die modellen wordt uiteindelijk bepaald door de
context en de kwaliteit van het onderzoek.


% ============================================================
\FVDsection{Samengevat}
% ============================================================

Samenhang betekent dat twee grootheden systematisch samen voorkomen of
samen veranderen.

Causaliteit betekent dat een verandering van de ene grootheid een
verandering van de andere veroorzaakt.

\[
\boxed{
\text{samenhang}
\neq
\text{causaliteit}
}
\]

Een samenhang kan onder andere beïnvloed worden door

\begin{itemize}
  \item een derde variabele;
  \item omgekeerde causaliteit;
  \item toeval;
  \item de manier waarop gegevens geselecteerd of verzameld zijn.
\end{itemize}

Observationeel onderzoek kan sterke samenhang aantonen, maar laat vaak
alternatieve verklaringen open.

Een goed ontworpen experiment kan sterkere causale conclusies
ondersteunen doordat de onderzochte factor doelbewust wordt veranderd
en andere relevante factoren zoveel mogelijk worden gecontroleerd.

Daarom geldt:

\[
\important{
\text{een hoge correlatie}
\not\Rightarrow
\text{een bewezen oorzaak-gevolgrelatie}
}
\]

en zelfs

\[
\important{
\text{een uitstekend regressiemodel}
\not\Rightarrow
\text{een causale verklaring}
}
\]

De juiste slotvraag is steeds:

\[
\boxed{
\text{Welke conclusie laten de gegevens én de onderzoeksopzet toe?}
}
\]

\end{document}
